\section*{Capítulo \ref{ch:b3:bacteriology} --- Bacteriología: crecimiento, fisiología y genética}

\begin{solution}{exo:b3:bacteriology:1}
\omterm{def:b3:bacteriology:envelope}{Grampositiva}: membrana citoplasmática envuelta en una pared gruesa de
\omterm{def:b3:bacteriology:envelope}{peptidoglicano} de muchas capas, atravesada por ácidos teicoicos.
\omterm{def:b3:bacteriology:envelope}{Gramnegativa}: una capa fina de \omterm{def:b3:bacteriology:envelope}{peptidoglicano} en un periplasma y después
una \omterm{def:b3:bacteriology:envelope}{membrana externa} cuya monocapa exterior es \omterm{def:b3:bacteriology:envelope}{lipopolisacárido}, con
porinas. El complejo de cristal violeta y yodo queda atrapado en la pared
gruesa y deshidratada por el alcohol de la \omterm{def:b3:bacteriology:envelope}{grampositiva} (morada), pero se
lava fuera de la \omterm{def:b3:bacteriology:envelope}{gramnegativa} a través de su pared fina y su \omterm{def:b3:bacteriology:envelope}{membrana
externa} alterada, que toma entonces la contratinción rosa.
\end{solution}

\begin{solution}{exo:b3:bacteriology:2}
$10^{6}$ veces $= 2^{19.9}$: unas $20$ generaciones; $T = \qty{8}{h}/20
= \qty{24}{min}$; $\mu = \ln 2/\qty{0.4}{h} = \qty{1.7}{h^{-1}}$.
\end{solution}

\begin{solution}{exo:b3:bacteriology:3}
\omterm{def:b3:bacteriology:hgt}{Transformación} --- Griffith (1928) y Avery, MacLeod y McCarty (1944):
los neumococos virulentos muertos por calor transforman a los avirulentos
vivos, y el agente es ADN. \omterm{def:b3:bacteriology:hgt}{Transducción} --- Zinder y Lederberg (1952): el
fago P22 lleva genes de \emph{Salmonella} de una cepa a otra a través de
un filtro que detiene a las células. \omterm{def:b3:bacteriology:hgt}{Conjugación} --- Lederberg y Tatum
(1946): dos cepas auxótrofas de \emph{E.~coli} dan recombinantes
prototrofas solo al mezclarse, y hace falta contacto celular (el tubo en
U de Davis).
\end{solution}

\begin{solution}{exo:b3:bacteriology:4}
Penicilina: las transpeptidasas que entrecruzan el \omterm{def:b3:bacteriology:envelope}{peptidoglicano}, que
las células animales no tienen. Estreptomicina: la subunidad ribosómica
pequeña, cuya forma bacteriana difiere de la eucariota. Ciprofloxacino:
la ADN-girasa, una topoisomerasa bacteriana ausente en el ser humano.
Rifampicina: la subunidad $\beta$ de la ARN-polimerasa bacteriana, que la
enzima eucariota no comparte.
\end{solution}

\begin{solution}{exo:b3:bacteriology:5}
$S^{*} = K_{S}D/(\mu_{\max} - D)$, $X^{*} = Y(S_{0} - S^{*})$. $D = 0.2$:
$S^{*} = \qty{0.0067}{g/L}$, $X^{*} = \qty{0.80}{g/L}$. $D = 0.6$:
$S^{*} = 0.06$, $X^{*} = 0.78$. $D = 0.79$: $S^{*} = 1.58$, $X^{*} = 0.17$.
Arrastre cuando $D = \mu(S_{0}) = 0.8\times 2/2.02 = \qty{0.79}{h^{-1}}$.
\end{solution}

\begin{solution}{exo:b3:bacteriology:6}
$(N/V)_{\text{umbr}} = kA_{\text{umbr}}/p$. Órgano luminoso: $0.01\times
6\times 10^{12}/100 = 6\times 10^{8}$ por litro $= 6\times 10^{5}$ por
mL. Agua de mar: $10\times 6\times 10^{12}/100 = 6\times 10^{11}$ por
litro $= 6\times 10^{8}$ por mL --- mil veces más, y nunca se alcanza en
el mar.
\end{solution}

\begin{solution}{exo:b3:bacteriology:7}
\qty{25}{mm}: sensible, CMI baja. \qty{12}{mm}: sensibilidad reducida,
intermedia --- puede fracasar a dosis corrientes. \qty{0}{mm}:
resistente, con crecimiento hasta el disco. El fármaco difunde desde el
disco de modo que su concentración cae con la distancia de forma
reproducible; el crecimiento se detiene donde la concentración iguala a
la CMI, así que el radio del halo y la CMI se relacionan por una curva de
calibración establecida con cepas de CMI conocida.
\end{solution}

\begin{solution}{exo:b3:bacteriology:8}
El \omterm{def:b3:genetic-engineering:tools}{plásmido} conjugativo lleva un integrón cuya serie de casetes contiene
varios genes de resistencia, más transposones que llevan más; el \omterm{def:b3:genetic-engineering:tools}{plásmido}
cruza en un solo apareamiento y los expresa todos. Se agrupan porque la
selección por cualquiera de los fármacos mantiene todo el \omterm{def:b3:genetic-engineering:tools}{plásmido}, de
modo que los genes que viajan de gorra en él persisten (la
coselección); los integrones capturan casetes; los transposones saltan a
los \omterm{def:b3:genetic-engineering:tools}{plásmidos}; y la unidad de transferencia es el \omterm{def:b3:genetic-engineering:tools}{plásmido}, no el gen.
\end{solution}

\begin{solution}{exo:b3:bacteriology:9}
Un solo fármaco: $10^{-7}\times 10^{9} = 100$ células resistentes
esperadas. Dos fármacos independientes:
$10^{-14}\times 10^{9} = 10^{-5}$. Las persistentes no son mutantes: una
pequeña fracción de células latentes sobrevive a cualquier fármaco y
vuelve a crecer como una población sensible cuando el fármaco se retira,
de modo que un segundo fármaco no ayuda; solo un tratamiento lo bastante
largo como para atraparlas al despertar, o un fármaco activo sobre
células latentes.
\end{solution}

\begin{solution}{exo:b3:bacteriology:10}
$X > X^{*}$: el consumo $\mu X/Y$ supera al suministro, de modo que $S$
cae por debajo de $S^{*}$; entonces $\mu(S) < D$ y
$\mathrm{d}X/\mathrm{d}t < 0$, y $X$ vuelve a bajar; al bajar $X$, $S$
sube de nuevo. Las dos desviaciones se corrigen: el estado estacionario
es estable. Dos especies: cada una necesita $\mu_{i}(S) = D$ para
persistir, lo que fija su propio $S_{i}^{*}$; el nutriente solo puede
tomar un valor, de modo que persiste a lo sumo una especie. La que tiene
el $S^{*}(D)$ más bajo lleva $S$ por debajo de lo que la otra necesita, y
la otra es arrastrada fuera; cuál sea esa especie puede cambiar con $D$
si sus curvas de Monod se cruzan.
\end{solution}

\begin{solution}{exo:b3:bacteriology:11}
Con la densidad $\rho = N/V$, el estado apagado $A = p_{0}\rho/k$ existe
mientras se mantenga por debajo de $A_{\text{umbr}}$, es decir,
$\rho < kA_{\text{umbr}}/p_{0}$; el estado encendido $A = p_{1}\rho/k$
existe mientras se mantenga por encima, es decir, $\rho
> kA_{\text{umbr}}/p_{1}$. Para $kA_{\text{umbr}}/p_{1} < \rho <
kA_{\text{umbr}}/p_{0}$ los dos son coherentes consigo mismos: una
población que se ha encendido se mantiene encendida con su producción
mayor, y una que no lo ha hecho sigue apagada. Utilidad: el compromiso
--- una vez que una colonia ha decidido construir una biopelícula o
liberar toxina, una bajada de densidad no deshace la decisión, y el
cambio es brusco y sincrónico en lugar de gradual.
\end{solution}

\begin{solution}{exo:b3:bacteriology:12}
(a) Dosis altas en tandas cortas: el nivel se mantiene por encima de la
CMI de los mutantes y mata por igual al tipo silvestre y a los mutantes
de primer paso --- a favor, con el límite de la toxicidad. (b) Dosis
bajas en tandas largas: el nivel se queda en la ventana durante días y
cría resistencia --- en contra. (c) Interrumpir cuando ceden los
síntomas: los supervivientes son las células menos sensibles y el nivel
que baja atraviesa la ventana --- en contra. (d) Combinación: una célula
ha de ser resistente a los dos, cosa que mil millones de células no
contienen --- a favor.
\end{solution}

\begin{solution}{pb:b3:bacteriology:1}
\textbf{1.} $18$ generaciones: $2^{18} = 2.6\times 10^{5}$ células,
$2.6\times 10^{3}$ por mL.
\textbf{2.} $2\times 10^{9}\times 100 = 2\times 10^{11}$ células $=
2^{37.5}$: al cabo de \qty{12.5}{h}; masa \qty{0.2}{g}.
\textbf{3.} $6\times 10^{27}/10^{-12} = 6\times 10^{39} = 2^{132}$
células: $132$ generaciones, \qty{44}{h}. Los nutrientes, el oxígeno y
los desechos lo detienen en un día; el crecimiento exponencial siempre es
breve.
\textbf{4.} Las células estacionarias han desmontado sus ribosomas y
reprimido sus enzimas metabólicas; han de reconstruirlos antes de
dividirse. Las células exponenciales llegan del todo equipadas.
\textbf{5.} $\mu = \ln 2/\qty{0.333}{h} = \qty{2.08}{h^{-1}}$; con $T =
\qty{60}{min}$, \qty{0.69}{h^{-1}}; tras \qty{6}{h}, $2^{6} = 64$ en
lugar de $2^{18}$: un factor $2^{12} \approx 4000$ menos.
\textbf{6.} $200/(0.1\times 10^{-6}) = 2\times 10^{9}$ células viables
por mL. El microscopio cuenta además las células muertas y las que no
crecerán en la placa, y un grumo da una sola colonia.
\textbf{7.} $D = 0.5/1 = \qty{0.5}{h^{-1}}$; $T = \ln 2/0.5 =
\qty{1.39}{h} = \qty{83}{min}$.
\textbf{8.} $S^{*} = 0.01\times 0.5/0.5 = \qty{0.01}{g/L}$; $X^{*} =
0.5\times 1.99 \approx \qty{1}{g/L} = 10^{12}$ células por litro,
$10^{9}$ por mL.
\textbf{9.} $0.5\times 10^{12} = 5\times 10^{11}$ células por hora; $720/1.39
\approx 520$ generaciones al mes.
\textbf{10.} $D = 0.95$: $S^{*} = 0.01\times 0.95/0.05 = \qty{0.19}{g/L}$,
$X^{*} = 0.5\times 1.81 = \qty{0.9}{g/L}$. Con $F = \qty{1.1}{L/h}$, $D
> \mu_{\max}$: arrastre.
\textbf{11.} El mutante mantiene $S^{*} = 0.01\times 0.5/0.7 = \qty{0.0071}{g/L}$;
a esa glucosa la cepa original crece a $1.0\times 0.0071/0.0171 =
\qty{0.42}{h^{-1}} < D$ y decae como $e^{-0.08t}$: queda arrastrada fuera
en unas pocas semanas. El mutante se adueña del recipiente.
\textbf{12.} $10^{-9}\times 10^{12} = 1000$ mutaciones beneficiosas por
generación: la adaptación es continua, con un clon ventajoso nuevo
barriendo la población cada pocos centenares de generaciones --- el
quimiostato es una máquina de evolución.
\textbf{13.} $kA_{\text{umbr}}/p = 0.05\times 6\times 10^{12}/500 =
6\times 10^{8}$ por litro $= 6\times 10^{5}$ células por mL.
\textbf{14.} $10^{10}/6\times 10^{5} \approx 1.7\times 10^{4}$ veces.
\textbf{15.} $20\times 6\times 10^{12}/500 = 2.4\times 10^{11}$ por litro,
$2.4\times 10^{8}$ por mL --- seis órdenes de magnitud por encima de las
$10^{2}$ por mL del mar: oscuras.
\textbf{16.} Diez veces son $3.3$ duplicaciones, \qty{6.6}{h}. A $10^{9}$
por mL la población sigue $1700$ veces por encima del umbral: la luz
sigue encendida.
\textbf{17.} En el mar nadie ve la luz y no alimenta a ningún
hospedador: un mutante oscuro se ahorra el \qty{20}{\%} y supera a los
que brillan. En el órgano el calamar selecciona a favor de la luz ---
expulsa o no alimenta a las tramposas oscuras ---, de modo que el coste
compra un hogar.
\textbf{18.} Bloquéese el receptor o destrúyase el \omterm{def:b3:bacteriology:quorum}{autoinductor} (el
apagado del quórum): las bacterias viven pero nunca despliegan sus
toxinas, y el sistema inmunitario las elimina. Como las bacterias
sensibles y las ``resistentes'' crecen igual --- no se mata a nadie ---,
la ventaja selectiva de escapar al fármaco es pequeña, y la resistencia
debería propagarse más despacio que frente a un fármaco que mata.
\textbf{19.} A: $10^{-8}\times 10^{9} = 10$; B: $100$; ambos: $10^{-15}\times
10^{9} = 10^{-6}$.
\textbf{20.} $P_{0}(\text{A}) = e^{-10} = 5\times 10^{-5}$;
$P_{0}(\text{ambos}) = e^{-10^{-6}} = 0.999999$.
\textbf{21.} $16 \to 8$: una semivida, \qty{4}{h}; $16 \to 1$: cuatro,
\qty{16}{h}; dentro de la ventana de las \qty{4}{h} a las \qty{16}{h},
\qty{12}{h} por dosis.
\textbf{22.} Con una dosis cada \qty{12}{h}: dentro de la ventana de la
hora $4$ a la $12$, dos tercios del tiempo; a lo largo de diez días los
mutantes de primer paso de A --- diez de ellos al principio --- quedan
seleccionados y A por sí solo fracasa.
\textbf{23.} Por encima de \qty{8}{mg/L} en todo momento: dósese cada
semivida (\qty{4}{h}), o bien una dosis que alcance un máximo de
\qty{64}{mg/L} cada \qty{12}{h} ($64 \to 8$ en tres semividas), o bien
una perfusión continua. Costes: la toxicidad de los máximos altos, o seis
dosis al día que los pacientes no cumplirán.
\textbf{24.} $10^{-4}\times 10^{9} = 10^{5}$ persistentes sobreviven pase
lo que pase con los fármacos; cuando cesa el tratamiento despiertan y
vuelven a crecer como una población del todo sensible --- una recaída. El
tratamiento ha de durar lo bastante como para atraparlas al salir de la
latencia, que es por lo que la tuberculosis dura seis meses.
\textbf{25.} $10^{9}$ células por mL en el quimiostato; la luz se
enciende a $6\times 10^{5}$ células por mL; y $P_{0} = 0.999999$ de que
la tanda de dos fármacos no encuentre ninguna célula resistente.
\end{solution}
